% SOURCE_AUTHORITY: sga5_fr_workpass.tex, lines 15266--15274 (LF-normalized unit).
% SOURCE_SHA256: 791F4EFFC5E02832D5D77ED03518C8156D6F07E4C8238B03545DB93D883FBB28
Demostración del lema 2: como $M=T(L)$ es proyectivo en cada grado, la hipótesis del lema se expresa diciendo que $T(\varphi)$ es homotópico a $\varphi_0$; existe, por tanto, un operador de homotopía $k_0=(k_0^i)_i$, con $k_0^i:M^i\to M^{i-1}$, tal que
\[
T(\varphi)-\varphi_0=d_0k_0+k_0d_0,
\]
donde $d_0$ es la diferencial de $M$. Por el lema 3(i), $k_0^i$ se levanta a una aplicación lineal $k^i:L^i\to L^{i-1}$. Se obtiene así un operador de homotopía $k=(k^i)_i$ en $L$. Basta poner
\[
\varphi'=\varphi-kd-dk
\]
($d=$ diferencial de $L$) para comprobar el lema 2.
